% authority: supporting
% status: draft/control
% task: RT-20260718-037
\section{Refuter stress of the cycle--boundary selector candidate}

\paragraph{Scope.}
This is one bounded Refuter stress of exact proposal-only
\(\mathrm{EqSrcCycleBoundarySelectorLaw}_{\rm src}^{\rm cand,v1}\), source
hash
\texttt{64901302bd6f37819ebd30aaaebaf39611d2c71c97d45b6cea6a39df11e0df94},
after the fresh audit with hash
\texttt{6169702f60659488e05cf0386afedc4b58b7395926923d5a2d365512182ca78f}.
It supplies no repair, adoption, canonical-ontology edit,
physical-admissibility claim, general \(\mathrm{EqSrc}\) discharge, ledger
change, or downstream promotion.

\subsection{Exact conditional core retained}

For an admitted oriented marked record \(S=(X,P)\), put \(Q=X\setminus P\),
\[
 p=\mathbf1_P,\qquad q=\mathbf1_Q,\qquad
 A_S=\operatorname{span}_{\mathbb F_2}\{p,q\},\qquad
 B_S=\mathbb F_2q.
\]
The candidate relation is congruence modulo \(B_S\), and its declared
variations are translations by \(0\) and \(q\).  Conditional on the supplied
oriented mark and on the candidate's admitted domain, the quotient
factorization, marked-bijection naturality, inverse and composition
certificates, and declared-translation invariance remain valid.  This packet
does not refute those conditional algebraic statements.

\subsection{Three-line equivariance obstruction}

\paragraph{Proposition (no unmarked equivariant line selector).}
Let \(A=\mathbb F_2^2\) be supplied without a distinguished basis, mark, or
line.  Its three nonzero proper subspaces are
\[
 L_p=\langle p\rangle,\qquad L_q=\langle q\rangle,\qquad
 L_d=\langle p+q\rangle .
\]
The action of \(\operatorname{GL}(2,2)\) on this set of lines is transitive.
Consequently no nonzero proper line is invariant under all unmarked linear
automorphisms, and there is no \(\operatorname{GL}(2,2)\)-equivariant
constant line selector on the bare two-dimensional presentation.

\paragraph{Proof.}
Every nonzero vector spans one of the three lines.  Given two nonzero
vectors, extend each to a basis and map the first basis to the second.  The
resulting invertible linear map carries the first line to the second.  Hence
the three lines form one orbit.  A line fixed by the whole group would have a
singleton orbit, which none has. \(\square\)

This is a presentation-level obstruction, not a claim that the physical
source gauge group is \(\operatorname{GL}(2,2)\).  It shows that the
distinguished line enters through the decorated mark or some additional
source structure; the bare \(A\)-space does not select it.

\paragraph{Corollary (equality--universal collapse after selector erasure).}
Let \(R_p,R_q,R_d\) be the congruence relations for the three lines.  Their
intersection is equality because
\(L_p\cap L_q\cap L_d=\{0\}\).  Their generated join is the universal
relation because \(L_p+L_q+L_d=A\).  Thus enforcing full line-permuting
symmetry by intersection or generated join yields the two degenerate
extremes, not a fourth nondegenerate line quotient.

\subsection{Minimal orientation-descent countermodel}

Take
\[
 X=\{0,1\},\qquad P=\{0\},\qquad Q=\{1\}.
\]
Both \(S_P=(X,P)\) and \(S_Q=(X,Q)\) satisfy the candidate's odd-mark
condition.  Under the common identification
\(A=\operatorname{span}\{p,q\}\),
\[
 B_{S_P}=\langle q\rangle,\qquad B_{S_Q}=\langle p\rangle.
\]
The first relation has classes
\(\{0,q\}\), \(\{p,p+q\}\); the second has classes
\(\{0,p\}\), \(\{q,p+q\}\).  They are distinct.

\paragraph{Proposition (presentation covariance without orientation
descent).}
The marked-source functor remains natural within its decorated groupoid, but
the candidate does not descend through the forgetful operation
\((X,P)\mapsto(X,\{P,Q\})\) on this even-carrier branch.

\paragraph{Proof.}
The two orientations have the same forgotten unordered bipartition and the
same state space under the identity identification, but their selected lines
and kernel-pair relations differ.  Descent would require the two outputs to
agree under that identification.  They do not. \(\square\)

This is not an inverse or cocycle defect inside the marked groupoid.
Identity, inverse, and composition there survive.  The failure is a descent
failure when orientation is forgotten.  It is therefore invalid to infer
that orientation is gauge, irrelevant, or safely forgettable.

The linear complement swap \(c(\alpha p+\beta q)=\beta p+\alpha q\)
transports \(B_{S_P}\) to \(B_{S_Q}\) and satisfies \(c^2=1\).  Hence weak
descent up to abstract pair isomorphism survives.  It retains only the
isomorphism class of a two-dimensional space with a line; it does not select
one relation on actual source states.  When \(|P|\ne|Q|\), no carrier
bijection exchanges the two blocks, so this algebraic swap need not arise
from a source transformation at all.

\paragraph{Domain correction.}
The exact candidate requires odd \(|P|\), not oddness of both blocks.
Therefore:
\begin{itemize}
  \item if \(|X|\) is even, \(Q\) is also odd and the complement orientation
  is admitted, so the descent countermodel applies;
  \item if \(|X|\) is odd, \(Q\) is even and only the odd block is admitted,
  so orientation can be reconstructed from the unordered bipartition on
  that restricted subdomain and complementing the candidate input changes a
  defined output into \(\Bottom\); and
  \item on a bare carrier without a partition, the prior full-symmetric-group
  obstruction to a nontrivial natural mark remains.
\end{itemize}
This corrects the scope of the stress target without editing the candidate or
the prior audit.

\subsection{Field and cardinality are sufficient, not necessary}

Let \(k\) be a field and retain the two-orbit fixed space
\(A_k=\operatorname{span}_k\{p,q\}\).  The candidate-style sum functional
satisfies
\[
 \ell_P(\alpha p+\beta q)=|P|\,\alpha\quad\hbox{in }k.
\]

\paragraph{Proposition (characteristic--cardinality branch theorem).}
If \(\operatorname{char}(k)\nmid |P|\), then
\(\ker(\ell_P|_{A_k})=\langle q\rangle\).  If \(k\) has positive
characteristic dividing \(|P|\), then
\(\ell_P|_{A_k}=0\), the kernel is all of \(A_k\), and the candidate's
proper-kernel test fails.

\paragraph{Proof.}
The displayed formula gives
\(\ell_P(p)=|P|1_k\) and \(\ell_P(q)=0\).  The coefficient of \(p\) is
therefore killed exactly when \(|P|1_k=0\). \(\square\)

Thus \(\mathbb F_2\) with odd \(|P|\) is one sufficient branch, not an
algebraically unique field/parity choice.  Moreover, because every vector in
\(A_k\) is constant on \(P\), evaluation of that constant,
\[
 e_P(\alpha p+\beta q)=\alpha ,
\]
is marked-natural and has kernel \(\langle q\rangle\) for every nonempty
\(P\), without a parity restriction.  This alternative is not adopted or
installed here.  It demonstrates that the candidate's field and parity
package is underselected by the desired kernel relation.

\subsection{Exact affine subgroup gap}

For one selected line \(B=\langle q\rangle\), the affine group
\(\operatorname{AGL}(2,2)\) has 24 elements.  The relation stabilizer is
\[
 G_{\rm rel}=\{z\mapsto Lz+t:L(B)=B\},
\]
the pointwise quotient-label fixer is
\[
 G_{\rm lab}=\{z\mapsto Lz+t:L(B)=B,\ t\in B\},
\]
and the declared variation group is
\[
 J=\{z\mapsto z+b:b\in B\}.
\]

\paragraph{Proposition (strict \(2<4<8\) variation hierarchy).}
\[
 |J|=2,\qquad |G_{\rm lab}|=4,\qquad |G_{\rm rel}|=8,
\]
with strict inclusions
\[
 J\cong C_2\;<\;G_{\rm lab}\cong V_4\;<\;
 G_{\rm rel}\cong D_8.
\]

\paragraph{Proof.}
The stabilizer of \(B\) in \(\operatorname{GL}(2,2)\) has two elements.
Arbitrary translation gives eight relation stabilizers.  Requiring the
quotient label itself to remain fixed restricts translation to the two
elements of \(B\), giving four maps; restricting the linear part to the
identity gives the two declared translations.  The order-four group is
elementary abelian.  The order-eight semidirect product contains an
order-four affine element and five involutions, hence is the dihedral group
of order eight. \(\square\)

More precisely, \(J\triangleleft G_{\rm lab}\triangleleft G_{\rm rel}\),
\(J=Z(G_{\rm rel})\), and both successive quotients have order two.

All three groups preserve the relation at their stated scope.  The algebra
therefore does not select \(J\) as the unique exhaustive or physical
variation class.  Exact kernel-aligned robustness is conditional
codetermination from the same mark, not an independent physical-selection
certificate.

\paragraph{Finite-field extension.}
For \(A=\mathbb F_s^2\), there are \(s+1\) boundary lines and
\[
\begin{aligned}
 |\operatorname{AGL}(2,s)|&=s^3(s-1)^2(s+1),\\
 |G_{\rm rel}|&=s^3(s-1)^2,\\
 |G_{\rm lab}|&=s^2(s-1),\\
 |J_{\rm full}|&=s,
\end{aligned}
\]
where \(J_{\rm full}\) contains all translations along the boundary line.
The \(s=2\) values are \(24,8,4,2\).  Carrying only the two original
translations to odd characteristic is not composition-closed; in larger
characteristic-two fields it is a proper order-two subgroup of the full
boundary translations.  This generalization reinforces, rather than
resolves, the missing field and physical-variation selection.

\subsection{Collapse, nonuniqueness, and surviving certificates}

The required failure branches are now explicit:
\begin{enumerate}
  \item \emph{collapse:} in positive characteristic dividing \(|P|\), the
  sum-functional kernel is all of \(A\), so the partial selector fails;
  \item \emph{nonuniqueness:} the bare two-dimensional space has three
  symmetry-related proper lines and no full-linear-symmetry invariant one;
  \item \emph{inverse/cocycle:} no defect occurs inside the exact decorated
  marked groupoid;
  \item \emph{orientation descent:} complementary admitted orientations on
  even carriers produce distinct relations after forgetting orientation;
  and
  \item \emph{variation fragility:} \(C_2<V_4<D_8\), with no independent
  current-ontology rule selecting the smallest class as physical.
\end{enumerate}

The surviving core is substantial but conditional: marked-source naturality,
quotient factorization and relation-level uniqueness, and robustness under
the two declared translations all pass.  What fails is the premise that
current ontology independently selects the oriented mark, its acquisition
provenance, the boundary line, the field/cardinality rule, the physical
morphism category, and the physical variation class.

\subsection{Refuter classification, freeze, and next route}

The exact Refuter classification is
\[
 \texttt{scoped\_obstruction}.
\]
The minimal countermodel is the two-point even-carrier complement pair above.
It blocks an orientation-independent, current-ontology-derived, independently
physical reading of this exact candidate.  It is not a global no-go theorem:
a richer source-side provenance law, an independently specified variation
category, an orientation-descent or irrelevance theorem on a justified
domain, or a distinct source extension may remain possible.

The exact
\(\mathrm{EqSrcCycleBoundarySelectorLaw}_{\rm src}^{\rm cand,v1}\)
construction--audit--stress cycle is now
\texttt{locally\_frozen}.  Repeating its construction, audit, or stress
without a materially different source law or theorem is barred.  General
\(\mathrm{EqSrc}\), the milestone, and future source extensions are not
frozen.

The next lawful route is one bounded
\texttt{theoretical-continuation-selector@0.1.0}
\texttt{ontology-law-research-packet}.  It must choose a materially distinct
source-provenance plus independent-variation witness, an orientation-descent
or irrelevance theorem, a distinct scoped no-go question, or local freeze.
It may not repeat, repair, adopt, or promote this exact candidate.

\paragraph{Distance-to-GR and forbidden conclusions.}
The source-equivalence milestone remains at \texttt{draft object exists};
general \(\mathrm{EqSrc}\), \(\mathrm{RetainH}\), \(\mathrm{GenH}\),
\(M_{\rm src}\), \(g_{\rm eff}\), matter coupling, Einstein equations, and
exact-GR recovery remain open.  Neither authoritative ledger changes.
No canonical ontology edit, source-law adoption, physical admissibility or
covariance, benchmark promotion, Gate Chair verdict, completed derivation,
future source-extension impossibility, program-level no-go, or global theory
rejection follows.
